\documentclass[11pt]{article}

\usepackage{amsmath,amssymb}
\usepackage{xurl}
\usepackage[hidelinks]{hyperref}
\setlength{\emergencystretch}{2em}

% Shared commands for notes in docs/paper-gaps/.
% Notation follows the LDT paper (references/ldt-paper/).

\newcommand{\Lean}{\textsc{Lean}}
\newcommand{\leanid}[1]{\textsf{#1}}
\newcommand{\ghissue}[1]{\href{https://github.com/LionSR/MIPStarRE/issues/#1}{GitHub issue \##1}}
\newcommand{\ghpr}[1]{\href{https://github.com/LionSR/MIPStarRE/pull/#1}{GitHub PR \##1}}

% --- Paper-compatible notation ---
\newcommand{\F}{\mathbb{F}}
\newcommand{\N}{\mathbb{N}}
\newcommand{\C}{\mathbb{C}}
\newcommand{\tr}{\operatorname{tr}}
\newcommand{\ot}{\otimes}
\newcommand{\E}{\mathop{\bf E\/}}
\newcommand{\bx}{\mathbf{x}}
\newcommand{\by}{\mathbf{y}}
\newcommand{\bu}{\mathbf{u}}
\newcommand{\bv}{\mathbf{v}}

% Approximation relation (paper uses \approx_\eps and \simeq_\zeta)
\newcommand{\approxeps}{\approx_{\varepsilon}}
\newcommand{\simgeq}{\simeq_{\zeta}}

% Polynomial spaces (paper uses \polyfunc, \polysub, \polymeas)
\newcommand{\polyfunc}[3]{\operatorname{Poly}(#1,#2,#3)}
\newcommand{\polysub}[3]{\operatorname{PolySub}(#1,#2,#3)}
\newcommand{\polymeas}[3]{\operatorname{PolyMeas}(#1,#2,#3)}


\title{Issue 1228: The SVD Step in Orthonormalization}
\author{}
\date{2026-05-05}

\begin{document}
\maketitle

\section{At a glance}

\begin{itemize}
  \item \textbf{Difficulty: medium.}  The paper's SVD step is finite-dimensional
  linear algebra, and the project already has the main positive-Gram route.
  \item \textbf{Estimated weight:} \textbf{3--6 theorem interfaces} and
  \textbf{300--700 lines} for the local route; a reusable rectangular complex
  SVD/polar theorem would likely be \textbf{800--1500 lines}.
  \item \textbf{Mathlib/project split:} about \textbf{65\% Mathlib} and
  \textbf{35\% project-local}.  Mathlib supplies spectral theory, adjoints,
  orthonormal bases, and square roots; the project supplies the rectangular
  $X,\widehat X,P$ packaging and estimates.
  \item \textbf{Key Mathlib inputs:} finite-dimensional inner-product spaces,
  matrix adjoints, Hermitian eigenvectors/eigenvalues, positive semidefinite
  matrices, and continuous-functional-calculus square roots.%
  \footnote{Representative APIs include \leanid{Orthonormal},
  \leanid{Matrix.IsHermitian.eigenvectorBasis},
  \leanid{Matrix.IsHermitian.eigenvalues}, \leanid{CFC.sqrt}, and
  \leanid{CFC.sqrt\_unique}.}
\end{itemize}

\section{Key theorem forms to develop}

\begin{enumerate}
  \item \textbf{Rectangular polar theorem.}  If $X$ is $m\times d$, $m\le d$,
  and $Q=X^\dagger X$, construct $\widehat X$ with
  \[
    \widehat X\widehat X^\dagger=I_m,
    \qquad
    X^\dagger\widehat X=\sqrt Q.
  \]
  \item \textbf{SVD adapter.}  If $X=U\Sigma V^\dagger$ and
  $\widehat X=UI_{m\times d}V^\dagger$, derive the two identities above from
  the unitary laws and positivity of the SVD middle factor.
  \item \textbf{Projective repair estimate.}  For
  $P_a=\widehat X^\dagger T_a\widehat X$, prove $P_aP_b=\mathbf 1_{a=b}P_a$ and
  the $P$-versus-$Q$ estimate from the squared-difference bound.
\end{enumerate}

\section{The source step}

The SVD appears in the orthonormalization proof of
\cite{Ji2020LowIndividualDegree}.%
\footnote{The SVD passage is in
\path{references/ldt-paper/orthonormalization.tex}, lines~774--845.}
After rounding and rank reduction, the paper has projectors $Q_a$ with
$Q=\sum_aQ_a$ and total rank at most $d$.  It defines
\[
  X_a=\sum_i\ket{a,i}\bra{v_{a,i}},
  \qquad
  X=\sum_aX_a,
  \tag{\texttt{[blueprint] def:matrix-decomposition-Q}}
\]
with $X_a=T_aX$ and $Q_a=X^\dagger T_aX$.  Since $X$ is $m\times d$ with
$m\le d$, the paper takes
\[
  X=U\Sigma_{m\times d}V^\dagger.
  \tag{\texttt{[blueprint] def:svd-of-X}}
\]

\section{What the SVD supplies}

The paper replaces the singular values by $1$:
\[
  \widehat X=UI_{m\times d}V^\dagger,
  \qquad
  P_a=\widehat X^\dagger T_a\widehat X.
  \tag{\texttt{[blueprint] def:projective-P}}
\]
It uses three consequences: \textbf{coisometry}
$\widehat X\widehat X^\dagger=I$, \textbf{square-root coupling}
$X^\dagger\widehat X=\sqrt Q$, and
\[
  (X-\widehat X)(X-\widehat X)^\dagger
  \le (XX^\dagger-I)^2.
\]
The last inequality is the scalar fact $(\sigma-1)^2\le(\sigma^2-1)^2$ for
$\sigma\ge0$.  These facts give
\[
  Q_a\ot I\approx_{30\zeta^{1/4}}P_a\ot I,
  \tag{\texttt{lem:P-Q-approx}}
\]
and hence $A_a\ot I\approx_{84\zeta^{1/4}}P_a\ot I$.%
\footnote{The consequences are in
\path{references/ldt-paper/orthonormalization.tex}, lines~984--1064, using the
auxiliary identities in lines~862--896.  The $P$-versus-$Q$ estimate and final
triangle step are in lines~1089--1192.}

\section{Blueprint and formalization}

The blueprint follows the same sequence: matrix decomposition, SVD, definition of
$P$, then the three identities and the $P$-versus-$Q$ estimate.%
\footnote{See \path{blueprint/src/chapter/ch04_projective.tex}, lines~704--785
and 833--1153.}
Lean separates \textbf{packaging} from \textbf{linear algebra}: the data package
stores the needed $X,\widehat X,P$ identities, while the substantive algebra is
provided either by explicit SVD adapters or by the positive-Gram construction.%
\footnote{The formal names include \leanid{svdOfX}, \leanid{QXPLayerData},
\leanid{rectangularSvd\_xHat\_coisometry},
\leanid{rectangularSvd\_xHat\_mixed\_raw},
\leanid{exists\_xHat\_of\_positive\_gram\_spectrum}, and
\leanid{pQApprox\_ofRankReductionSigmaRangePositiveGram}; see
\doclink{MIPStarRE/LDT/MakingMeasurementsProjective/QXPLayer/Core.lean} and
\doclink{MIPStarRE/LDT/MakingMeasurementsProjective/QXPLayerIdentities.lean}.}
The version of Mathlib pinned here lacks exactly the rectangular polar theorem
above, so the project-local positive-Gram route supplies it.%
\footnote{The local coisometry helper
\leanid{Matrix.exists\_mul\_conjTranspose\_eq\_one\_of\_card\_le} is in
\doclink{MIPStarRE/Quantum/FiniteHilbert.lean}.  Possible upstreaming is tracked
in \ghissue{1250}.}

\section{Conclusion}

There is \textbf{no SVD discrepancy}.  The paper needs $\widehat X$ with
$\widehat X\widehat X^\dagger=I$ and $X^\dagger\widehat X=\sqrt Q$, plus the
squared-difference estimate.  The formalization preserves this mechanism via
explicit SVD adapters or the positive-Gram route.

\bibliographystyle{alpha}
\bibliography{references}

\end{document}
